% ID: SEL-D0042-OB-E04
% ESTADO: cotejada. Segunda pasada visual y borrador compilado revisados.
% Fuente: Selectividad/sel_2014_matematicas_aplicadas/EXAMEN -ANDAL-1.pdf
% SHA-256: 375424f4e0d3c0409619ebaae4ae4d20c441a2585efe3dc74eb98d765114311e
% Página: 2; opción B; ejercicio original 4.

Queremos estudiar la proporción de personas de una población que usan una determinada marca de ropa; para ello se hace una encuesta a 950 personas y se obtiene que 215 de ellas usan esa marca. Utilizando un contraste de hipótesis $(H_0:p\geq0.25)$:
\begin{enumerate}[label=\alph*)]
\item \textbf{(1.5 puntos)} \textquestiondown{}Podemos afirmar con estos datos y con un nivel de significación del 5\% que al menos el 25\% de toda la población usa esa marca de ropa?
\item \textbf{(1 punto)} \textquestiondown{}Y con un nivel de significación del 1\%?
\end{enumerate}
